\section{Conclusion}
\label{sec:ch4_conclusion}

This chapter has introduced \textbf{OmniFORE}, a novel framework for optimizing resource forecasts in edge-cloud networks. By integrating the components of the OmniFORE framework, including attention-based time-series modeling, temporal clustering, strategic data sampling, and generalization techniques, OmniFORE generalizes across workloads and efficiently predicts diverse workloads in volatile environments.

Our approach uses strategic data sampling to extract representative traces, employs temporal clustering to differentiate workload patterns, and uses attention mechanisms to capture both short-term and long-term dependencies. Extensive experiments on real-world datasets demonstrated that \ac{OmniFORE} consistently outperforms existing benchmark schemes across various scenarios, including regular and zero-shot predictions. The framework showed significant improvements in prediction accuracy, inference speed, and adaptability to both low and high variance data.

These contributions address the complexities of cloud-native workloads and enable more adaptive edge-cloud systems. However, while \ac{OmniFORE} addresses the prediction and generalization challenges, generating accurate forecasts is only half the battle. The final challenge lies in acting upon these predictions to make real-time orchestration decisions in complex, distributed environments. This necessitates a transition from predictive intelligence to autonomous decision-making.

The next chapter addresses this need by introducing \textbf{AgentEdge}, a distributed multi-agent framework designed to close the loop by autonomously executing orchestration actions based on the insights generated by predictive models like \ac{AERO} and \ac{OmniFORE}.
